\documentclass{article}
% General document formatting
\usepackage[utf8]{inputenc}
% Related to math
\usepackage{amsmath,amssymb,amsfonts,amsthm}
\pagenumbering{gobble}
\usepackage{geometry}
\geometry{margin=1.5in}
%\usepackage{fullpage}
\title{Research Plan Joint PhD KULeuven - VUB}
\author{Emilio Gamba\\ Promotors: \textit{Tias Guns (KULeuven) and Bart Bogaerts (VUB)}}
\date{June 2022}

\begin{document}
	
	\makeatletter
	\renewcommand\paragraph{\@startsection{paragraph}{4}{\z@}%
		{1.25ex \@plus1ex \@minus.2ex}
		{-0.5em}%
		{\normalfont\normalsize\bfseries}}
	\makeatother
	
	\maketitle
	
	This document gives an overview of the research activities that I’m planning for the remainder of my PhD. 
	The research activities reflect my current interest in how to explain constraint satisfaction and optimisation problems. 
	In the first activity, I will investigate how to efficiently compute explanations under uncertainty. 
	The second activity is related to optimizing explanation sequences for constraint satisfaction problems.
	Finally, the last activity will be devoted to writing the PhD dissertation and preparing for the oral defence.
	
	\section{Research Plan}
	The research activities in this research plan built upon the already published papers where I am the main author \cite{ijcai2021-191} or a contributor  \cite{bogaerts2020step,bogaerts2021framework}.
	
	\subsection{Research activity 1: \emph{Explaining constraint solving under uncertainty}}
	\paragraph{Role}Main author
	\paragraph{Target Period} July 2022 – February 2023
	\paragraph{Target}IJCAI23 (conference paper)
	\paragraph{Collaborations}Maxime Mulamba (VUB-KULeuven), Bart Bogaerts (VUB), Tias Guns (KULeuven)
	
	\paragraph{Context} In a real-life context, many combinatorial problems, such as Sudoku puzzles, come under the form of an image. Such problems require the integration of perception, using, for instance Machine Learning (ML), and a reasoning task to solve a combinatorial problem from visual input. 
	Using the techniques for integrating ML and Constraint Solving developed by PhD researcher Maxime Mulambda, we investigate how to provide explanations where the input is uncertain, for instance probabilistic estimates.
	To do this we build on the cost-optimal explanation techniques we have developed so far. 
	
%	\newpage
	\subsection{Research activity 2: \emph{Optimizing explanation sequences}}
	\paragraph{Role}Main author
	\paragraph{Target Period}January 2023 – August 2023
	\paragraph{Target}IJCAI/CP-23 (conference paper)
	\paragraph{Collaborations}Bart Bogaerts (VUB), Tias Guns (KULeuven)
	
	\paragraph{Context} In previous work, Tias Guns, Bart Bogaerts and myself have studied how to generate a sequence of \emph{simple} explanations, where difficulty of an explanation is quantified using a cost function. The focus of that work has been on generating these explanations efficiently and guarantee that they are optimal with respect to the cost-function. There explanation steps are generated independently.
	While the explanations are `simple', there is no guarantee that the sequence as a whole is simple as well.
	
	In this research activity, we provide better \emph{sequences} of explanations.
	For instance, a sequence is easier if there is a relation between subsequent steps to avoid constant context switches which make the reasoning harder to understand for humans. 
	Therefore, in order to improve the generated sequences, we investigate alternative measures of complexity of an explanation sequence.
	The focus  here is on the algorithmic challenges that arise from improving the sequence and not the aspects related to cognitive complexity. 
	The goal is to develop generic (with respect to the complexity measure used) global algorithms building on previous work to provide an optimal explanation sequence.
	
	%% Own focus on algorithms not cognitive
	
	%\subsection{Research activity 3: \emph{Abstraction – Refinement of explanation sequences}}
	%\paragraph{Role}Main author
	%\paragraph{Target Period} April 2023 – September 2023
	%\paragraph{Target}AAAI 24 (conference paper)
	%\paragraph{Collaborations}Bart Bogaerts (VUB), Tias Guns (KULeuven)
	%
	%
	%\paragraph{Context} The explanation sequences generated to explain how to solve constraint satisfaction problems are highly dependent on the way the problem is modelled and encoded into the solver.
	%Using the complexity measures developed for optimizing explanation sequences in Research Activity 2, we investigate how to zoom-in or abstract parts of the explanation sequence. 
	
	\subsection{Research activity 3: \emph{Writing the PhD dissertation \& Preparation for the defence} }
	\paragraph{Target Period} September 2023- March 2024
	\vspace*{1cm} \\
	June 2022, Brussels, Belgium,\\
	Emilio Gamba (PhD Researcher) \hspace{1.5in} Tias Guns (Promotor)\\
	\vspace*{1cm} \\
	\bibliographystyle{plain}
	\bibliography{reference}
	
\end{document}